\section{Recommendations}
\label{sec:recs}

These are addressed to benchmark maintainers and to anyone reporting agent
evaluations, ourselves included. They are cheap, and most of them are fields
rather than experiments.

\begin{enumerate}\itemsep2pt
\item \textbf{Record the serving endpoint, not just the model identifier.}
  Provider and quantization at minimum. Two entries carrying the same model id
  today can be different weights on different hardware, and a reader has no way
  to tell. This is the single highest-value change here, and it is one column.
\item \textbf{Record image digest, task-set commit and run date.} Digest rather
  than tag: a tag moves. The run date is what lets a reader notice that two
  entries are separated by an unannounced change in endpoint behaviour.
\item \textbf{Report the number of runs, and a range when there is more than
  one.} Three runs with a minimum and maximum costs three times the compute and
  removes most of the ambiguity. Where a board already stores a run count,
  surfacing it in the rendered table is free.
\item \textbf{State a minimum meaningful gap in the leaderboard
  documentation.} The same-system difference threshold of
  Section~\ref{sec:noise} is one way to compute one; resampling whole rounds
  puts it between 5.18 and 8.53 points. Differences below that range between
  single runs are not resolvable. Saying so protects submitters as much as
  readers: it removes the incentive to claim a lead that a re-run would
  erase.
\item \textbf{Log completion tokens per action.} It is the cheapest available
  witness to an endpoint changing underneath a run. Aggregate token counts are
  dominated by the prompt and will not show it.
\item \textbf{Distinguish budget exhaustion from task failure in the recorded
  outcome.} A run stopped at the round cap and a run that did the wrong thing
  are different events; one in five published runs is the former.
\item \textbf{Three small fixes we can offer upstream}, each independent of
  everything above: merge rather than assign caller-supplied request options;
  guard the response-content accessor against a \texttt{None} body, which a
  reasoning model returns when it exhausts its budget mid-thought; and match
  container cleanup against the configured image rather than a hardcoded tag.
\end{enumerate}

There is a stronger conclusion available, and our data supports it: if a
serving configuration is not part of an API contract, then an evaluation built
on hosted inference is not reproducible in principle, only in practice and only
until the provider changes something. The strict remedy is to self-host and
freeze the weights. We are not proposing that everyone do so --- for most
groups the cost is prohibitive and the hosted endpoint is the thing they
actually deploy against --- but a field that cites two-point differences should
be explicit that it has chosen convenience over reproducibility, rather than
assuming it has both.

None of this makes a benchmark deterministic, and none of it should. The
purpose is narrower: to let a reader of a leaderboard tell the difference
between a system that improved and an instrument that moved.
